\documentclass[a4paper,11pt]{article}
\usepackage[utf8]{inputenc}
\usepackage[T1]{fontenc}
\usepackage[ngerman]{babel}
\usepackage{amsmath,amssymb}
\usepackage{xcolor}
\usepackage{colortbl}
\usepackage{array}
\usepackage{tikz}
\usepackage{enumitem}
\usepackage[margin=2cm]{geometry}
\definecolor{bgdark}{HTML}{1E1E1E}
\definecolor{fgtext}{HTML}{E6E6E6}
\definecolor{gridcol}{HTML}{4A4A4A}
\definecolor{accent}{HTML}{6CB6FF}
\definecolor{loescol}{HTML}{7FD18C}
\pagecolor{bgdark}
\color{fgtext}
\arrayrulecolor{fgtext}
\setlength{\parindent}{0pt}
\newcommand{\karo}[1]{\par\noindent\begin{tikzpicture}\draw[step=5mm,gridcol,very thin](0,0) grid (\linewidth,#1);\end{tikzpicture}\par\vspace{2mm}}
\newcommand{\aufgabe}[2]{\vspace{3mm}\par\noindent{\large\color{accent}\textbf{Aufgabe #1}}\hfill{\color{gridcol}\normalsize (#2 Punkte)}\par\vspace{1mm}}

\begin{document}

% ====================== KOPF ======================
\begin{center}
  {\LARGE\color{accent}\textbf{Optimierungsverfahren}}\\[3pt]
  {\large Probeklausur --- Teil 2: Optimierung}\\[2pt]
  {\small M. Zaefferer (Stil) \quad$\bullet$\quad Übungsklausur, klausurnah}
\end{center}
\vspace{2mm}
\noindent\textbf{Name:}\ \underline{\hspace{6cm}}\hfill\textbf{Matrikelnr.:}\ \underline{\hspace{4cm}}\\[4pt]
\noindent\textbf{Datum:}\ \underline{\hspace{4cm}}\hfill\textbf{Bearbeitungszeit:}\ ca.\ 60 Minuten\\[4pt]
\noindent\textbf{Gesamtpunkte:}\ 50 \hfill \textbf{Hilfsmittel:}\ Open-Book (eigene Unterlagen/Formelsammlung), Taschenrechner\\[4pt]
{\color{gridcol}\rule{\linewidth}{0.4pt}}
\vspace{2mm}

\noindent\textbf{Punktetabelle} (wird von der Korrektur ausgefüllt):
\par\vspace{1mm}
\renewcommand{\arraystretch}{1.5}
\noindent\begin{tabular}{|>{\centering\arraybackslash}m{1.6cm}|c|c|c|c|c|c|c|c|c||c|}
\hline
\textbf{Aufgabe} & 2.1 & 2.2 & 2.3 & 2.4 & 2.5 & 2.6 & 2.7 & 2.8 & 2.9 & $\Sigma$ \\\hline
\textbf{max} & 6 & 5 & 7 & 5 & 6 & 6 & 4 & 6 & 5 & 50 \\\hline
\textbf{erreicht} & & & & & & & & & & \\\hline
\end{tabular}
\renewcommand{\arraystretch}{1}
\par\vspace{3mm}
\noindent{\small\color{gridcol}Hinweise: Schreiben Sie Ihren Lösungsweg nachvollziehbar auf. Reine Ergebnisse ohne Rechenweg geben i.\,d.\,R. keine volle Punktzahl.}

% ====================== AUFGABE 2.1 ======================
\clearpage
\aufgabe{2.1 --- Modellierung}{6}
Gegeben ist die folgende Problemstellung:
\begin{itemize}[nosep,leftmargin=1.4em]
  \item Ein Unternehmen betreibt zwei Online-Werbekanäle: \emph{Influencer-Marketing} (IM) und \emph{Display-Werbung} (DW).
  \item Für beide Kanäle steht ein gemeinsames Werbebudget von $12\,000\,€$ zur Verfügung.
  \item Die Anzahl der Neukunden $z_1$ aus IM wächst mit dem eingesetzten Budgetanteil $b_1$ etwa wie $z_1=\sqrt{b_1}$, die Neukunden $z_2$ aus DW wie $z_2 = 0.5\cdot\sqrt{b_2}$.
  \item Der Gewinn pro IM-Kunde beträgt $40\,€$, pro DW-Kunde $90\,€$.
  \item Es soll der erwartete Gewinn maximiert werden.
\end{itemize}
\vspace{1mm}
\textbf{(a.)} (4 Punkte) Definieren Sie klar die Entscheidungsvariablen und formulieren Sie das mathematische Optimierungsproblem mit Zielfunktion und Nebenbedingungen. Eine Lösung ist \emph{nicht} notwendig.
\par\vspace{1mm}
\textbf{(b.)} (2 Punkte) Handelt es sich um ein lineares Optimierungsproblem? Begründen Sie kurz.
\karo{12cm}

% ====================== AUFGABE 2.2 ======================
\clearpage
\aufgabe{2.2 --- Grafische Lösung \& Slackvariablen}{5}
Gegeben ist das folgende \textbf{Minimierungsproblem}:
\[
\min_{\vec x}\ f(\vec x)=x_1+x_2 \quad\text{u.d.N.}\quad
\begin{aligned}[t]
 & g_1(\vec x): 3x_1+2x_2 \ge 12\\
 & g_2(\vec x): x_1+2x_2 \ge 8\\
 & g_3(\vec x): x_1 \le 10\\
 & g_4(\vec x): x_2 \le 9\\
 & g_5(\vec x): x_1+x_2 \le 20\\
 & x_1\ge 0,\ x_2\ge 0
\end{aligned}
\]
In der Grafik sind die aktiven Nebenbedingungen $g_1$ (durchgezogen) und $g_2$ (durchgezogen), der zulässige Bereich (grau) und zwei Isolinien der Zielfunktion (blau gestrichelt) dargestellt; $g_3,g_4,g_5$ liegen weit außerhalb des gezeichneten Ausschnitts.
\par\vspace{2mm}
\begin{center}
\begin{tikzpicture}[scale=0.62]
  \fill[white,opacity=0.06] (2,3)--(0,6)--(0,10)--(10,10)--(10,0)--(8,0)--(2,3);
  \draw[gridcol,very thin,step=1] (0,0) grid (10,10);
  \draw[->,fgtext] (0,0)--(10.5,0) node[right]{$x_1$};
  \draw[->,fgtext] (0,0)--(0,10.5) node[above]{$x_2$};
  \foreach \i in {2,4,6,8,10}{\node[fgtext,below] at (\i,0){\small\i}; \node[fgtext,left] at (0,\i){\small\i};}
  % g1: 3x1+2x2=12 -> (4,0),(0,6)
  \draw[accent,very thick] (4,0)--(0,6) node[above right,pos=0.15]{\small$g_1$};
  % g2: x1+2x2=8 -> (8,0),(0,4)
  \draw[accent,very thick] (8,0)--(0,4) node[above right,pos=0.3]{\small$g_2$};
  % isolines x1+x2=5 and =9
  \draw[blue!60!white,dashed,thick] (5,0)--(0,5);
  \draw[blue!60!white,dashed,thick] (9,0)--(0,9);
  \node[blue!60!white] at (4.3,4.3){\small$f=9$};
  \node[blue!60!white] at (2.0,1.0){\small$f=5$};
\end{tikzpicture}
\end{center}
\textbf{(a.)} (3 Punkte) Lösen Sie das Problem mit Hilfe der Grafik. Geben Sie $\vec x^{\,*}$ und den optimalen Zielfunktionswert $f(\vec x^{\,*})$ an.
\par\vspace{1mm}
\textbf{(b.)} (2 Punkte) Welche Slackvariablen der Nebenbedingungen $g_1$ bis $g_5$ nehmen für Ihre Lösung einen Wert ungleich Null an?
\karo{6.5cm}

% ====================== AUFGABE 2.3 ======================
\clearpage
\aufgabe{2.3 --- Transportproblem}{7}
Gegeben sei folgendes klassisches Transportproblem: Von $3$ Produktionsstätten sollen täglich Produkte zu $4$ Läden transportiert werden. Die Produktionsstätten produzieren täglich: $200,\ 180,\ 120$ Produkte. Die Läden sollen täglich verkaufen: $120,\ 130,\ 150,\ 100$.
Die Transportkosten (Zeilen $=$ Produktionsstätten, Spalten $=$ Läden) sind:
\[
C=\begin{pmatrix} 4 & 8 & 8 & 5\\ 6 & 5 & 9 & 7\\ 9 & 7 & 3 & 2 \end{pmatrix}
\]
Gegeben ist zudem folgender Lösungsvorschlag:
\[
X=\begin{pmatrix} 120 & 80 & 0 & 0\\ 0 & 50 & 130 & 0\\ 0 & 0 & 20 & 70 \end{pmatrix}
\]
\textbf{(a.)} (2 Punkte) Wie groß wären die Gesamttransportkosten für diesen Lösungsvorschlag?
\par\vspace{1mm}
\textbf{(b.)} (2 Punkte) Der Lösungsvorschlag ist keine zulässige Lösung. Warum?
\par\vspace{1mm}
\textbf{(c.)} (3 Punkte) Ändern Sie \emph{eine} Spalte des Lösungsvorschlags, sodass dieser zulässig wird.
\karo{10cm}

% ====================== AUFGABE 2.4 ======================
\clearpage
\aufgabe{2.4 --- Grafische nichtlineare Optimierung}{5}
Gegeben ist eine Zeichnung für das folgende nichtlineare \textbf{Minimierungsproblem}:
\[
\min_{\vec x}\ f(\vec x)=x_1\cdot\sqrt{x_2}
\quad\text{mit}\quad
\begin{aligned}[t]
 & h(\vec x): x_1^2-x_2=0\\
 & g_1(\vec x): x_2>0,\qquad g_2(\vec x): -2\le x_1\le 2
\end{aligned}
\]
Die zulässige Menge ist die rot gezeichnete Kurve (Parabel $x_2=x_1^2$) im Bereich $-2\le x_1\le 2$.
\par\vspace{2mm}
\begin{center}
\begin{tikzpicture}[scale=1.05]
  \draw[gridcol,very thin,step=0.5] (-2.4,0) grid (2.4,4.4);
  \draw[->,fgtext] (-2.5,0)--(2.6,0) node[right]{$x_1$};
  \draw[->,fgtext] (0,0)--(0,4.6) node[above]{$x_2$};
  \foreach \i in {-2,-1,1,2}{\node[fgtext,below] at (\i,0){\small\i};}
  \foreach \j in {1,2,3,4}{\node[fgtext,left] at (0,\j){\small\j};}
  % Parabel als zulaessige Kurve
  \draw[red!75!white,very thick,smooth] plot coordinates
    {(-2,4)(-1.5,2.25)(-1,1)(-0.5,0.25)(0,0)(0.5,0.25)(1,1)(1.5,2.25)(2,4)};
  \node[red!75!white] at (1.65,3.5){\small$h$};
  % Grenzen x1 = -2 und x1 = 2
  \draw[accent,dashed] (-2,0)--(-2,4.4) node[above]{\small$x_1=-2$};
  \draw[accent,dashed] (2,0)--(2,4.4) node[above]{\small$x_1=2$};
\end{tikzpicture}
\end{center}
Bestimmen Sie das globale Minimum mit Hilfe der Zeichnung. Geben Sie sowohl die Koordinaten $\vec x$ als auch den Zielfunktionswert $f(\vec x)$ an.
\karo{6cm}

% ====================== AUFGABE 2.5 ======================
\clearpage
\aufgabe{2.5 --- KKT-Bedingungen}{6}
Stellen Sie für das folgende nichtlineare Optimierungsproblem die KKT-Bedingungen auf (die Gleichungen / Ungleichungen, anhand derer das Problem gelöst werden könnte). \emph{Sie müssen das Problem nicht lösen.}
\[
\min_{\vec x}\ f(\vec x)=(x_1-2)^2+x_2^2
\]
mit
\[
h(\vec x):\ x_1+x_2-1=0,\qquad g(\vec x):\ x_1^2+x_2^2\le 4
\]
\karo{12cm}

% ====================== AUFGABE 2.6 ======================
\clearpage
\aufgabe{2.6 --- Nelder-Mead}{6}
Gegeben ist ein Minimierungsproblem mit Zielfunktion $f(\vec x)=x_1^2+x_2^2$ und zwei Variablen $\vec x\in\mathbb{R}^2$.
Gegeben sind zudem die folgenden $3$ Punkte (bereits sortiert), die einen Simplex im Sinne des Nelder-Mead-Algorithmus bilden:
\[
\vec x_0=\begin{pmatrix}1\\2\end{pmatrix},\quad
\vec x_1=\begin{pmatrix}3\\2\end{pmatrix},\quad
\vec x_2=\begin{pmatrix}4\\4\end{pmatrix}
\]
\[
f(\vec x_0)=5,\qquad f(\vec x_1)=13,\qquad f(\vec x_2)=32
\]
Berechnen Sie mit diesen Angaben und Parameter $\alpha=1$ den Reflektionsschritt und bestimmen Sie, welcher Schritt auf die Reflektion folgen muss.
\karo{10cm}

% ====================== AUFGABE 2.7 ======================
\clearpage
\aufgabe{2.7 --- $(1+1)$-Evolutionsstrategie}{4}
Nehmen Sie an, die $(1+1)$-ES wird für ein Minimierungsproblem eingesetzt. Die Speicherlänge für die Erfolgsratenberechnung beträgt $20$. In den letzten $19$ Schritten wurden $5$ Erfolge erzielt. Die Qualität der aktuellen Lösung ist $f(x)=0.812$, die Qualität der mutierten Lösung ist $f(x^{\text{next}})=0.798$. Die aktuelle Schrittweite beträgt $\sigma=0.5$, der Schrittweitenmultiplikator ist $a=1.2$.
\par\vspace{1mm}
Bestimmen Sie die Schrittweite für den nächsten Schritt.
\karo{9cm}

% ====================== AUFGABE 2.8 ======================
\clearpage
\aufgabe{2.8 --- Non-Dominated-Sorting-Rang}{6}
Gegeben sind die Koordinaten von sieben Punkten im Zielraum, mit drei Zielfunktionen:
\begin{center}
\begin{tabular}{llllllll}
$P_1=(1,1,1)$ & $P_2=(2,2,1)$ & $P_3=(1,2,2)$ & $P_4=(2,2,2)$\\
$P_5=(2,0,2)$ & $P_6=(0,2,1)$ & $P_7=(1,0,3)$
\end{tabular}
\end{center}
Alle drei Zielfunktionen sollen \textbf{minimiert} werden. Geben Sie für jeden der Punkte den Non-Dominated-Sorting-Rang an.
\karo{10cm}

% ====================== AUFGABE 2.9 ======================
\clearpage
\aufgabe{2.9 --- Crowding Distance}{5}
Bestimmen Sie für die Punkte in der folgenden grafischen Darstellung des Zielraums die Crowding Distance (bei $2$ Zielfunktionen, Minimierung). Welche \textbf{vier} Punkte würde man anhand der Crowding Distance in einem evolutionären Algorithmus für die Mehrzieloptimierung wählen? Alle Koordinaten sind ganzzahlig.
\par\vspace{2mm}
\begin{center}
\begin{tikzpicture}[scale=0.7]
  \draw[gridcol,very thin,step=1] (0,0) grid (7,9);
  \draw[->,fgtext] (0,0)--(7.5,0) node[right]{$f_1$};
  \draw[->,fgtext] (0,0)--(0,9.5) node[above]{$f_2$};
  \foreach \i in {1,2,3,4,5,6}{\node[fgtext,below] at (\i,0){\small\i};}
  \foreach \j in {1,...,8}{\node[fgtext,left] at (0,\j){\small\j};}
  \foreach \p/\x/\y in {1/0/8,2/6/1,3/1/5,4/2/4,5/4/2}{
    \fill[accent] (\x,\y) circle (3.2pt);
    \node[fgtext,above right] at (\x,\y){\small$P_{\p}$};
  }
\end{tikzpicture}
\end{center}
\textit{Punkte:}\ $P_1=(0,8),\ P_2=(6,1),\ P_3=(1,5),\ P_4=(2,4),\ P_5=(4,2)$.
\karo{6.5cm}

\vfill
\begin{center}{\color{gridcol}\small--- Ende der Probeklausur (Teil 2) ---}\end{center}

\end{document}
